%% kapa: exam=1 round=27 year=2016 period=1 form=A subject=civil_law subject_ko=민법 seq=15 number=15
%% source: Q-Net 감정평가사 기출문제 (cst00309) artlSeq=5169402 file=extracted/1008/5169402/제27회 감정평가사 1차 문제/제27회 감정평가사 1교시문제(A형).pdf
%% conversion: kordoc 4.24.51 -> tools/gen1.py (auto); review=proofread-sonnet; flags=none
\begin{question}{15}
토지거래허가를 받지 않아 토지매매계약이 유동적 무효의 상태에 있는 경우에 관한 설명으로 옳지 않은 것은? (다툼이 있으면 판례에 따름)
\begin{choices}
  \item 위 매매계약이 확정적으로 무효로 됨에 있어서 귀책사유가 있는 자는 그 계약의 무효를 주장할 수 없다.
  \item 허가구역 지정이 해제되면 위 매매계약은 확정적 유효로 된다.
  \item 허가구역 지정기간이 만료되었음에도 허가구역 재지정을 하지 아니한 경우, 위 매매계약은 확정적 유효로 된다.
  \item 허가를 받으면 위 매매계약은 소급해서 유효로 되므로 허가 후에 새로 매매계약을 체결할 필요는 없다.
  \item 사기에 의하여 위 매매계약이 체결된 경우, 취소권자는 토지거래허가를 신청하기 전에 사기에 의한 계약의 취소를 주장하여 거래허가신청협력에 거절의사를 일방적으로 명백히 함으로써, 그 계약을 확정적으로 무효화 시킬 수 있다.
\end{choices}
\end{question}
\answer{1}
